\section{The Method}

Every result in this paper is an in-silico computation on the substrate, never a measurement against nature. The protocol is fixed and the same for all experiments, and it is designed so that a result is the gap between a run and a control rather than a number read off in isolation. The whole comparison rests on a reproducible suite \cite{cluesurf2026repo}.

The base rule is deterministic, with no randomness anywhere. It is a reversible permutation of the tones. Initial conditions are deterministic patterns or the vacuum, never random seeds. Scale is varied by changing the lattice size, not by averaging over draws. A run is bit-identical on two machines. This matters for the claims. A deterministic base that still yields the Born rule and the Tsirelson bound is a stronger result than one that feeds in quantum randomness by hand.

Every claim is paired with a control that could fail. If the result were an artifact, the control would expose it. Most often it is the lossy, erasing rule run against the conserving one, or a flat substrate run against the curved one. A result counts only as the difference between the run and its control. Where the conserving rule gives an exact zero, the lossy control gives a nonzero number equal to the charge it destroys.

Thresholds are fixed before running. Pass conditions are set in advance, and where possible they are exact integer or ratio equalities with no tolerance, so a result is an identity rather than a fit. The discipline borrowed from Chronoflux is one substrate-parameter assignment that must fit every arena at once, with the pass bars set before the data are seen.

Each experiment carries a stable code of the form E-XXX-NNNN, registered once in a catalog. The three letters mark the arena, GRV for gravity, RLT for relativity, SPN for spin, FLD for fluid, CSM for cosmology, MTH for method, FND for foundations, and so on. The four digits are the registry number. Each result is also given a depth grade. L3 is an emergent and novel result from the base rule with a control. L2 is known physics reproduced on the substrate. L1 is known math confirmed. L0 is a consistency note, kept and labeled as circular rather than dropped.

The evidence base is 615 experiments across 18 categories, 84 at L3, 389 at L2, 128 at L1, and 14 at L0, with 412 backing a specific claim in the papers. The anchors of this paper are \texttt{E-FND-0030}, \texttt{E-GRV-0039}, \texttt{E-GRV-0042}, \texttt{E-GRV-0040}, \texttt{E-GRV-0041}, and \texttt{E-MTH-0001}, alongside the ladder-push results \texttt{E-RLT-0039}, \texttt{E-CSM-0043}, \texttt{E-FLD-0009}, and \texttt{E-FLD-0010}.

Negative results are reported as findings. A push that turns out not to be cleanly buildable is recorded as such. Two negatives from this round are reported here. A metric from the bare tone correlations does not build, because the ternary tone field is short-range and carries no extended structure to reconstruct a metric from. A clean geometry-dependent settling time does not build, because the observable is too noisy. These count as results.
